
% bismi ALLAH ARRAHMAN ARRAHIM

\documentclass[11pt,a4paper]{article}

\usepackage[utf8]{inputenc}
\usepackage[T1]{fontenc}
\usepackage[french]{babel}
\usepackage{amsmath,amssymb,amsfonts,amsthm,mathtools}
\usepackage{booktabs}
\usepackage{geometry}
\usepackage{enumitem}
\usepackage{microtype}
\usepackage{hyperref}
\hypersetup{hidelinks}
\usepackage{tikz}
\usetikzlibrary{arrows.meta,positioning}

\geometry{margin=2.5cm}

% -------------------------------------------------
% Environnements
% -------------------------------------------------
\newtheorem{definition}{Définition}[section]
\newtheorem{assumption}{Hypothèse}[section]
\newtheorem{proposition}{Proposition}[section]
\newtheorem{theorem}{Théorème}[section]
\newtheorem{corollary}{Corollaire}[section]
\newtheorem{remark}{Remarque}[section]
\newtheorem{lemma}{Lemme}[section]
\newtheorem{example}{Exemple}[section]

% -------------------------------------------------
% Opérateurs
% -------------------------------------------------
\DeclareMathOperator{\prox}{prox}
\DeclareMathOperator{\argmin}{argmin}
\DeclareMathOperator{\argmax}{argmax}
\DeclareMathOperator{\dist}{dist}
\DeclareMathOperator{\conv}{conv}

% -------------------------------------------------
% Raccourcis
% -------------------------------------------------
\newcommand{\R}{\mathbb R}
\newcommand{\Mcal}{\mathcal M}
\newcommand{\Ccal}{\mathcal C}
\newcommand{\Pcal}{\mathcal P}
\newcommand{\Tcal}{\mathcal T}
\newcommand{\Rcal}{\mathcal R}
\newcommand{\Acal}{\mathcal A}
\newcommand{\Bcal}{\mathcal B}

\title{Routage variationnel PCH-ESM\\
\large Sélection de contexte comme enveloppe inférieure certifiée}
\author{}
\date{}

\begin{document}

\maketitle

\begin{abstract}
Ce manuscrit formalise le routage dans le cadre PCH en séparant strictement le noyau de purification et la décision de contexte. Le noyau PCH est supposé fixé ; le routage n'est pas une couche ajoutée, mais la comparaison de valeurs réduites issues des problèmes locaux. On obtient ainsi un problème d'optimisation sur \(\R^d\times\Ccal\), puis sa relaxation entropique sur le simplexe des contextes. La quantité centrale est la valeur réduite normalisée du problème local. Les outils de Moreau, Berge, Clarke et Rockafellar--Wets sont rappelés là où ils sont réellement utilisés, tandis que l'apport propre du manuscrit est la lecture du routage comme géométrie de décision induite par une famille finie de fonctions valeur. Une compatibilité explicite avec la décomposition spectrale et la spécialisation certifiée de PCH-ESM est également donnée.
\end{abstract}

\tableofcontents

% =================================================
\section{Introduction et motivation}
% =================================================

On considère un noyau PCH déjà défini et validé. Pour chaque entrée latente \(w\in\R^d\), il existe plusieurs contextes admissibles \(c\in\Ccal\), chacun muni d'un ensemble convexe admissible, fermé et non vide \(\Mcal_c\subset\R^d\), et d'une fonctionnelle locale propre.

Le point de départ est simple : dans PCH, la question \emph{quel contexte purifier ?} n'est pas extérieure au modèle. Elle est déjà contenue dans la famille des valeurs locales associées aux contextes. Le routage n'ajoute donc pas une couche artificielle ; il explicite la comparaison des branches d'énergie que PCH fournit déjà.

La raison de cette formulation est claire. Dans les approches de type \emph{mixture of experts}, la décision de routage est souvent portée par un score appris ou par un gate séparé. Ici, on cherche au contraire une décision \emph{induite} par la structure variationnelle locale elle-même. Le contexte retenu n'est pas produit par un module externe ; il est l'indice qui minimise une valeur optimale locale.

Cette lecture est aussi la plus compatible avec l'architecture PCH-ESM : la décomposition spectrale y produit une famille finie de branches ou d'espaces spécialisés, et le routage peut alors choisir la branche certifiée de plus faible valeur réduite. Le routage devient ainsi la continuité naturelle de la théorie PCH.

La contribution du manuscrit est la suivante :
\begin{itemize}[leftmargin=2em]
    \item le noyau PCH est conservé sans modification ;
    \item le routage est formulé comme une minimisation imbriquée sur \(\R^d\times\Ccal\) ;
    \item la sélection de contexte est gouvernée par la valeur réduite locale ;
    \item la relaxation entropique fournit un routage souple explicite ;
    \item la géométrie des cellules et des frontières est déduite de la famille finie des valeurs réduites ;
    \item la compatibilité avec PCH-ESM est rendue explicite via une famille finie de branches certifiées.
\end{itemize}

\begin{remark}
Le point original du manuscrit est visible dès l'introduction : le routage n'est pas décrit comme un mécanisme heuristique, mais comme une géométrie de décision induite par une famille de fonctions valeur.
\end{remark}

\begin{figure}[ht]
\centering
\begin{tikzpicture}[node distance=1.7cm, >=Latex, every node/.style={align=center}]
\node (w) [draw, rounded corners, minimum width=2.8cm, minimum height=0.9cm] {$w\in\R^d$};
\node (locals) [draw, rounded corners, right=2.7cm of w, minimum width=4.1cm, minimum height=1.05cm] {Problèmes locaux\\$\widetilde J_c(\cdot;w)$};
\node (values) [draw, rounded corners, right=2.7cm of locals, minimum width=3.8cm, minimum height=1.05cm] {Valeurs réduites\\$\widetilde V_c(w)$};
\node (argmin) [draw, rounded corners, below=1.25cm of values, minimum width=3.8cm, minimum height=1.05cm] {Sélection\\$\argmin_{c\in\Ccal}\widetilde V_c(w)$};
\node (chosen) [draw, rounded corners, left=2.7cm of argmin, minimum width=3.8cm, minimum height=1.05cm] {Contexte choisi\\$c^\star(w)$};

\draw[->] (w) -- (locals);
\draw[->] (locals) -- (values);
\draw[->] (values) -- (argmin);
\draw[->] (argmin) -- (chosen);
\end{tikzpicture}
\caption{Schéma minimal du routage variationnel : le contexte choisi est l'indice qui minimise la valeur réduite.}
\end{figure}

\section*{Table des notations principales}
\begin{center}
\begin{tabular}{@{}ll@{}}
\toprule
Notation & Signification \\ \midrule
\(\widetilde J_c(u;w)\) & énergie locale normalisée du contexte \(c\) \\
\(\widetilde V_c(w)\) & valeur réduite associée au contexte \(c\) \\
\(u_c^\star(w)\) & minimiseur local unique du problème normalisé \\
\(F(w)\) & valeur globale de routage, \(F(w)=\min_c \widetilde V_c(w)\) \\
\(\Acal(w)\) & ensemble des contextes actifs \\
\(\Omega_c^{\mathrm{str}}\) & cellule stricte du contexte \(c\) \\
\(\Bcal\) & ensemble des points de bord \\
\(\Rcal_\beta\) & fonctionnelle de routage souple \\
\(\Delta_K\) & simplexe des contextes \\ \bottomrule
\end{tabular}
\end{center}

% =================================================
\section{Problème local PCH et normalisation canonique}
% =================================================

\begin{assumption}[Espace latent]
On fixe un espace euclidien réel \(\R^d\). Les états latents purifiés vivent dans \(\R^d\).
\end{assumption}

\begin{assumption}[Famille de contextes]
On considère un ensemble fini de contextes admissibles
\[
\Ccal=\{1,\dots,K\}.
\]
À chaque contexte \(c\in\Ccal\) est associé un ensemble convexe non vide, fermé, \(\Mcal_c\subset\R^d\).
\end{assumption}

\begin{assumption}[Paramètres PCH]
Pour chaque contexte \(c\), on fixe des paramètres \(\lambda_{\mathrm{ctx},c}>0\), \(\lambda_{\mathrm{ridge},c}\ge 0\) et un paramètre proximal \(\eta_c>0\).
\end{assumption}

\begin{definition}[Énergie contextuelle]
Pour un contexte \(c\in\Ccal\), on définit l'énergie
\[
J_c(u;w)=\frac12\|u-w\|^2+\lambda_{\mathrm{ctx},c}\,\dist(u,\Mcal_c)^2+\lambda_{\mathrm{ridge},c}\,\|u\|^2.
\]
\end{definition}

\begin{definition}[Opérateur PCH associé]
Pour un contexte \(c\in\Ccal\), on définit
\[
\Pcal_c(w):=\argmin_{u\in\R^d} J_c(u;w).
\]
\end{definition}

\begin{proposition}[Convexité forte et unicité]
Pour tout contexte \(c\) et toute entrée \(w\), la fonctionnelle \(u\mapsto J_c(u;w)\) est \(\mu_c\)-fortement convexe, où
\[
\mu_c:=1+2\lambda_{\mathrm{ridge},c}.
\]
En particulier, \(\Pcal_c(w)\) est un singleton.
\end{proposition}

\begin{proof}
Le terme \(u\mapsto \frac12\|u-w\|^2+\lambda_{\mathrm{ridge},c}\|u\|^2\) est \((1+2\lambda_{\mathrm{ridge},c})\)-fortement convexe. La fonction \(u\mapsto \dist(u,\Mcal_c)^2\) est convexe parce que \(\Mcal_c\) est non vide, fermé et convexe. La somme d'une fonction fortement convexe et de fonctions convexes est fortement convexe, d'où l'existence et l'unicité du minimiseur.
\end{proof}

\begin{remark}
Le noyau PCH n'est pas modifié par ce document. On ne change ni la forme de \(J_c\), ni la structure des ensembles \(\Mcal_c\), ni la purification finale. Le routage est ajouté comme problème de sélection séparé.
\end{remark}

\begin{definition}[Fonction de base normalisée]
On introduit
\[
f_c(u):=\lambda_{\mathrm{ctx},c}\,\dist(u,\Mcal_c)^2+\lambda_{\mathrm{ridge},c}\,\|u\|^2,
\qquad
g_c(u):=\frac{1}{\mu_c}f_c(u).
\]
Alors
\[
\widetilde J_c(u;w)=\frac{1}{\mu_c}J_c(u;w)
=
g_c(u)+\frac{1}{2\mu_c}\|u-w\|^2.
\]
\end{definition}

\begin{proposition}[Justification de la normalisation]
Le critère de routage est défini par les valeurs normalisées, et non par les coûts bruts. Pour chaque contexte \(c\), la normalisation par \(\mu_c\) a les effets suivants :
\begin{enumerate}[leftmargin=2em]
\item elle ne change pas le minimiseur de \(J_c(\cdot;w)\) ;
\item elle compare les branches sur une même échelle liée à leur rigidité quadratique ;
\item elle transforme le problème local en problème proximal pour la fonction convexe \(g_c\) ;
\item elle rend explicite la curvure locale de la branche via le paramètre \(\mu_c\).
\end{enumerate}
\end{proposition}

\begin{proof}
Le point (1) est immédiat car la multiplication par un scalaire strictement positif ne change pas l'ensemble des minimisateurs. Le point (3) suit de l'identité
\[
\widetilde J_c(u;w)=g_c(u)+\frac{1}{2\mu_c}\|u-w\|^2,
\]
qui est la forme standard d'une enveloppe de Moreau de paramètre \(\mu_c\). Les points (2) et (4) formalisent le choix du critère de routage : on compare les branches sur une même échelle, de sorte que le ridge n'écrase pas artificiellement la décision, tout en gardant la géométrie de chaque branche visible dans la réduction.
\end{proof}

\begin{definition}[Valeur réduite du contexte]
Pour tout contexte \(c\) et toute entrée \(w\), on définit la valeur réduite
\[
\widetilde V_c(w):=\inf_{u\in\R^d}\widetilde J_c(u;w).
\]
Par unicité du minimiseur, on note ce point
\[
u_c^\star(w):=\argmin_{u\in\R^d}\widetilde J_c(u;w).
\]
\end{definition}

\begin{definition}[Enveloppe de Moreau associée au contexte]
On note l'enveloppe de Moreau de \(g_c\) au paramètre \(\mu_c\) par
\[
e_{\mu_c g_c}(w):=\inf_{u\in\R^d}\left\{g_c(u)+\frac{1}{2\mu_c}\|u-w\|^2\right\}.
\]
\end{definition}

\begin{proposition}[Enveloppe de Moreau et gradient]
Pour tout contexte \(c\), la fonction \(g_c\) est propre, convexe et semi-continue inférieurement. On a
\[
u_c^\star(w)=\prox_{\mu_c g_c}(w),
\qquad
\widetilde V_c(w)=e_{\mu_c g_c}(w).
\]
De plus, \(e_{\mu_c g_c}\) est de classe \(C^1\) sur \(\R^d\) et
\[
\nabla e_{\mu_c g_c}(w)=\frac{1}{\mu_c}\bigl(w-\prox_{\mu_c g_c}(w)\bigr)
=\frac{1}{\mu_c}\bigl(w-u_c^\star(w)\bigr).
\]
Par conséquent,
\[
\nabla \widetilde V_c(w)=\frac{1}{\mu_c}\bigl(w-u_c^\star(w)\bigr).
\]
\end{proposition}

\begin{proof}
La convexité et la semi-continuité inférieure de \(g_c\) découlent de celles de \(\dist(\cdot,\Mcal_c)^2\) et de \(\|\cdot\|^2\). Le résultat classique sur les enveloppes de Moreau d'une fonction propre, convexe et s.c.i. implique que \(e_{\mu_c g_c}\) est de classe \(C^1\) et que
\[
\nabla e_{\mu_c g_c}(w)=\frac{1}{\mu_c}\bigl(w-\prox_{\mu_c g_c}(w)\bigr).
\]
Comme \(\widetilde J_c(u;w)=g_c(u)+\frac{1}{2\mu_c}\|u-w\|^2\), les minimisateurs coïncident et la formule annoncée suit.
\end{proof}

\begin{remark}
Ces identités sont des résultats standards de l'analyse convexe. Elles sont rappelées ici uniquement là où elles servent le raisonnement \cite{Moreau1965,Rockafellar1970,BauschkeCombettes2017,RockafellarWets1998}.
\end{remark}

% =================================================
\section{Routage exact, géométrie des cellules et régularité de Clarke}
% =================================================

\begin{definition}[Fonctionnelle de routage]
Pour une entrée \(w\in\R^d\), on introduit la fonctionnelle de routage
\[
\mathcal E_{\mathrm{route}}(w;u,c):=\widetilde J_c(u;w),
\qquad
(u,c)\in \R^d\times\Ccal.
\]
\end{definition}

\begin{theorem}[Lower envelope, routage exact et géométrie de décision]
Pour tout \(w\in\R^d\), on a
\[
F(w):=\min_{c\in\Ccal}\widetilde V_c(w)
=
\inf_{(u,c)\in \R^d\times\Ccal}\mathcal E_{\mathrm{route}}(w;u,c).
\]
De plus, le contexte optimal vérifie
\[
\Acal(w):=\argmin_{c\in\Ccal}\widetilde V_c(w)
\]
et tout minimiseur global \((\hat u(w),\hat c(w))\) satisfait
\[
\hat u(w)=u_{\hat c(w)}^\star(w).
\]
La fonction \(F\) est continue et localement lipschitzienne. Elle est directionnellement différentiable, et pour tout \(h\in\R^d\),
\[
F'(w;h)=\min_{c\in\Acal(w)}\langle \nabla \widetilde V_c(w),h\rangle.
\]
Enfin, son sous-différentiel de Clarke est donné par
\[
\partial^C F(w)=\conv\{\nabla \widetilde V_c(w): c\in\Acal(w)\}.
\]
Sur toute cellule stricte \(\Omega_c^{\mathrm{str}}\), la fonction \(F\) est de classe \(C^1\) et
\[
\nabla F(w)=\nabla \widetilde V_c(w).
\]
\end{theorem}

\begin{proof}
Comme \(\Ccal\) est fini,
\[
\inf_{(u,c)\in \R^d\times\Ccal}\mathcal E_{\mathrm{route}}(w;u,c)
=
\min_{c\in\Ccal}\inf_{u\in\R^d}\widetilde J_c(u;w)
=
\min_{c\in\Ccal}\widetilde V_c(w).
\]
Le minimiseur en \(u\) est unique pour chaque contexte \(c\), ce qui donne \(\hat u(w)=u_{\hat c(w)}^\star(w)\). La continuité et la Lipschitzité locale de \(F\) résultent du fait que \(F\) est le minimum fini de fonctions \(C^1\). La formule directionnelle et le sous-différentiel de Clarke sont les règles standards du minimum fini de fonctions \(C^1\) localement lipschitziennes \cite{Clarke1983}. Enfin, sur une cellule stricte, le contexte actif est unique et stable par continuité ; la fonction \(F\) y coïncide donc avec une seule branche \(C^1\).
\end{proof}

\begin{proposition}[Correspondance active, Berge et frontières]
La correspondance active \(\Acal:\R^d\rightrightarrows\Ccal\) a des valeurs non vides et finies. Elle est semi-continue supérieurement. En particulier, sur toute cellule de dominance stricte, elle est localement constante. Ce point est la version finie du théorème du maximum de Berge \cite{Berge1963}.
\end{proposition}

\begin{proof}
Pour chaque \(c\),
\[
\{w\in\R^d:\ c\in\Acal(w)\}
=
\bigcap_{j\in\Ccal}\{w\in\R^d:\widetilde V_c(w)\le \widetilde V_j(w)\},
\]
intersection finie d'ensembles fermés, donc ensemble fermé. Comme \(\Ccal\) est fini, la correspondance est à valeurs finies non vides et fermées ; il s'agit donc d'un cas discret du maximum theorem de Berge. Sur une cellule stricte, les inégalités sont strictes et stables par continuité, d'où la constance locale.
\end{proof}

\begin{definition}[Cellules de contexte]
On définit la cellule stricte associée à un contexte \(c\in\Ccal\) par
\[
\Omega_c^{\mathrm{str}}:=\left\{w\in\R^d:\ \widetilde V_c(w)<\widetilde V_j(w)\ \text{pour tout }j\in\Ccal,\ j\neq c\right\}.
\]
On définit également la cellule de dominance faible
\[
\Omega_c:=\left\{w\in\R^d:\ c\in\argmin_{j\in\Ccal}\widetilde V_j(w)\right\}
=
\bigcap_{j\neq c}\{w\in\R^d:\widetilde V_c(w)\le \widetilde V_j(w)\}.
\]
Enfin, on note
\[
\Bcal:=\{w\in\R^d:\ \#\Acal(w)\ge 2\}
\]
l'ensemble des points de bord.
\end{definition}

\begin{proposition}[Partition de l'espace latent et frontières]
Pour chaque \(c\in\Ccal\), l'ensemble \(\Omega_c^{\mathrm{str}}\) est ouvert et les ensembles \(\Omega_c\) sont fermés. De plus, les cellules strictes sont deux à deux disjointes, et
\[
\R^d=\left(\bigsqcup_{c\in\Ccal}\Omega_c^{\mathrm{str}}\right)\sqcup \Bcal.
\]
Pour chaque paire \(c\neq j\), en posant
\[
\Gamma_{c,j}(w):=\widetilde V_c(w)-\widetilde V_j(w),
\qquad
\Sigma_{c,j}:=\{w\in\R^d:\Gamma_{c,j}(w)=0\},
\]
\(\Sigma_{c,j}\) est fermé ; si \(\nabla\Gamma_{c,j}(w_0)\neq 0\), alors \(\Sigma_{c,j}\) est localement une hypersurface \(C^1\) au voisinage de \(w_0\).
\end{proposition}

\begin{proof}
Chaque \(\widetilde V_c-\widetilde V_j\) est continue. Donc \(\Omega_c^{\mathrm{str}}\), défini par des inégalités strictes de fonctions continues, est ouvert, et \(\Omega_c\), défini par des inégalités larges, est fermé. Les cellules strictes sont disjointes par contradiction immédiate.

Tout point \(w\) a un minimiseur unique, auquel cas il appartient à une cellule stricte, ou plusieurs minimiseurs, auquel cas il appartient à \(\Bcal\). Cela donne la décomposition annoncée.

La fermeture de \(\Sigma_{c,j}\) découle de la continuité de \(\Gamma_{c,j}\). Si \(\nabla\Gamma_{c,j}(w_0)\neq 0\), le théorème des fonctions implicites s'applique à l'équation \(\Gamma_{c,j}(w)=0\), ce qui fournit une hypersurface \(C^1\) locale.
\end{proof}

\begin{corollary}[Stabilité sous marge stricte]
Supposons qu'en un point \(w\in\R^d\), le minimiseur de \(c\mapsto \widetilde V_c(w)\) soit unique, noté \(c^\star\), et qu'il existe une marge
\[
\delta(w):=\min_{c\neq c^\star}\bigl(\widetilde V_c(w)-\widetilde V_{c^\star}(w)\bigr)>0.
\]
Alors il existe un voisinage \(U\) de \(w\) tel que tout \(z\in U\) admet le même contexte routé \(c^\star\).
\end{corollary}

\begin{proof}
Par continuité finie des fonctions \(\widetilde V_c\), on choisit \(U\) tel que \(|\widetilde V_c(z)-\widetilde V_c(w)|<\delta(w)/3\) pour tous \(c\) et tout \(z\in U\). L'argument standard de comparaison des branches donne alors la stabilité locale du contexte optimal.
\end{proof}

% =================================================
\section{Réalisation algorithmique et profil de calcul}
% =================================================

Le routage exact peut être coûteux, car il exige la résolution complète de chaque problème contextuel. On introduit donc un certificat de calcul par pré-roulage court. Cette partie n'est pas un appendice : elle fournit un estimateur quantitatif de la valeur réduite et une stabilité calculable du choix de contexte.

\begin{proposition}[Profil de calcul]
Si le coût d'une résolution locale du contexte \(c\) à précision \(\varepsilon\) est noté \(T_c(\varepsilon)\), alors un routage exact par comparaison exhaustive coûte
\[
C_{\mathrm{exact}}(\varepsilon)=\sum_{c\in\Ccal} T_c(\varepsilon).
\]
Un routage appris par gate a un coût d'inférence approximatif
\[
C_{\mathrm{gate}}\approx T_{\mathrm{gate}}+T_{c^\star},
\]
mais exige une phase d'apprentissage distincte et n'offre pas, en lui-même, de certificat de marge sur le contexte choisi. Le routage variationnel exact est donc plus coûteux au moment de l'inférence, mais il est certifiable, naturellement compatible avec PCH, et ne requiert aucun gate séparé.
\end{proposition}

\begin{proof}
Le premier terme suit de la comparaison explicite de tous les contextes. Le second profil reflète le schéma classique gate-plus-expert : un gate est évalué puis un expert est exécuté. La différence conceptuelle est que le gate appris doit être entraîné et peut se tromper hors distribution, alors que le routage variationnel compare directement les valeurs locales et peut être certifié par marge.
\end{proof}

\begin{definition}[Pré-roulage court]
Pour un contexte \(c\), on pose
\[
u_c^{(0)}(w)=w,
\qquad
u_c^{(n+1)}(w)=\Tcal_{c,\eta}(u_c^{(n)}(w);w),
\qquad n\ge 0.
\]
\end{definition}

\begin{definition}[Valeur de pré-roulage]
Pour un entier \(N\ge 1\), on définit
\[
\widetilde V_c^{(N)}(w):=\widetilde J_c(u_c^{(N)}(w);w).
\]
\end{definition}

\begin{assumption}[Hypothèse de contraction]
Pour un contexte \(c\) fixé, on suppose que l'opérateur
\[
\Tcal_{c,\eta}(z;w):=\argmin_{u\in\R^d}\left\{\widetilde J_c(u;w)+\frac{1}{2\eta}\|u-z\|^2\right\}
\]
est contractant sur \(\R^d\) avec facteur \(\rho_c\in(0,1)\), pour tout \(w\) dans le domaine considéré.
\end{assumption}

\begin{proposition}[Convergence géométrique et approximation de la valeur réduite]
Sous l'hypothèse de contraction ci-dessus, la suite \(u_c^{(N)}(w)\) converge vers \(u_c^\star(w)\) et
\[
\|u_c^{(N)}(w)-u_c^\star(w)\|\le \rho_c^N\|u_c^{(0)}(w)-u_c^\star(w)\|.
\]
De plus, si \(\widetilde J_c(\cdot;w)\) est Lipschitz sur un compact contenant l'orbite \(\{u_c^{(N)}(w)\}_{N\ge 0}\cup\{u_c^\star(w)\}\), alors il existe une constante \(L_{c,w}<\infty\) telle que
\[
\bigl|\widetilde V_c^{(N)}(w)-\widetilde V_c(w)\bigr|\le L_{c,w}\rho_c^N\|u_c^{(0)}(w)-u_c^\star(w)\|.
\]
\end{proposition}

\begin{proof}
La convergence géométrique des itérés est la conclusion standard du théorème du point fixe de Banach pour une contraction de facteur \(\rho_c\). Comme \(\widetilde J_c\) est continue et \(C^1\), elle est Lipschitz sur tout compact. En évaluant la différence des valeurs sur le compact contenant l'orbite et le point fixe, on obtient l'inégalité annoncée.
\end{proof}

\begin{theorem}[Stabilisation du routage approché]
Supposons que le minimiseur de \(c\mapsto \widetilde V_c(w)\) soit unique, noté \(c^\star(w)\). Alors il existe \(N_0\ge 1\) tel que, pour tout \(N\ge N_0\),
\[
\hat c_N(w)=c^\star(w),
\qquad
\hat c_N(w)\in\argmin_{c\in\Ccal}\widetilde V_c^{(N)}(w).
\]
\end{theorem}

\begin{proof}
Comme l'ensemble \(\Ccal\) est fini, \(\argmin_{c\in\Ccal}\widetilde V_c^{(N)}(w)\) est non vide pour tout \(N\). En outre, pour chaque \(c\), \(\widetilde V_c^{(N)}(w)\to \widetilde V_c(w)\). La convergence est donc uniforme sur \(\Ccal\). Posons
\[
\delta:=\min_{c\neq c^\star(w)}\bigl(\widetilde V_c(w)-\widetilde V_{c^\star(w)}(w)\bigr)>0.
\]
Il existe alors \(N_0\ge 1\) tel que, pour tout \(N\ge N_0\) et tout \(c\in\Ccal\),
\[
\bigl|\widetilde V_c^{(N)}(w)-\widetilde V_c(w)\bigr|<\frac{\delta}{3}.
\]
Les mêmes comparaisons que dans la stabilité sous marge stricte donnent alors \(\hat c_N(w)=c^\star(w)\) pour tout \(N\ge N_0\).
\end{proof}

\begin{corollary}[Routage approché par certificat court]
On peut définir un routage approché par
\[
\hat c_N(w)\in\argmin_{c\in\Ccal}\widetilde V_c^{(N)}(w).
\]
Lorsque \(N\) est suffisamment grand, \(\hat c_N(w)\) coïncide avec le routage exact sous l'hypothèse d'unicité du minimiseur.
\end{corollary}

\begin{remark}
Le pré-routage court n'est pas un bloc isolé : il fournit un estimateur calculable de la valeur réduite et sa stabilité provient directement de la marge entre branches. Il relie ainsi la théorie à une procédure calculable.
\end{remark}

\begin{quote}
\small
Algorithme de routage certifié

Entrée : \(w\), ensemble fini \(\Ccal\), tolérance \(\varepsilon\).

1. Pour chaque contexte \(c\in\Ccal\), résoudre le problème local et calculer une approximation certifiée \(\widetilde V_c^{(N)}(w)\).

2. Choisir \(c^\star\in\argmin_{c\in\Ccal}\widetilde V_c^{(N)}(w)\).

3. Si une marge \(\delta>0\) est certifiée, arrêter dès que
\[
\bigl|\widetilde V_c^{(N)}(w)-\widetilde V_{c^\star}^{(N)}(w)\bigr|<\delta/3
\qquad\text{pour tout }c\neq c^\star.
\]

Sortie : contexte routé \(c^\star\) et minimiseur local \(u_{c^\star}^\star(w)\).
\end{quote}

\begin{remark}
La procédure précédente n'est pas une heuristique dissimulée : elle n'utilise que des valeurs locales, une marge certifiée, et un critère d'arrêt explicite.
\end{remark}

% =================================================
\section{Routage souple comme relaxation entropique}
% =================================================

Le routage dur est discret. Pour obtenir une relaxation convexe, on passe au simplexe des contextes. Ce passage est une relaxation explicite du problème discret.

\begin{definition}[Simplexe des contextes]
On note
\[
\Delta_K:=\left\{p\in\R^K:\ p_c\ge 0,\ \sum_{c=1}^K p_c=1\right\},
\qquad
\Delta_K^\circ:=\{p\in\Delta_K:\ p_c>0\ \text{pour tout }c\}.
\]
\end{definition}

\begin{definition}[Fonctionnelle de routage souple]
Pour \(\beta>0\), on définit
\[
\Rcal_\beta(w,p)
=
\sum_{c=1}^K p_c\,\widetilde V_c(w)
+\frac{1}{\beta}\sum_{c=1}^K p_c\log p_c,
\qquad p\in\Delta_K,
\]
avec la convention \(0\log 0=0\).
\end{definition}

\begin{proposition}[Existence et unicité du routage souple]
Pour tout \(w\), la fonctionnelle \(p\mapsto \Rcal_\beta(w,p)\) admet un unique minimiseur \(p^\star_\beta(w)\) dans \(\Delta_K\).
\end{proposition}

\begin{proof}
La fonction \(p\mapsto \sum_c p_c\log p_c\) est strictement convexe sur \(\Delta_K\) et \(p\mapsto \sum_c p_c\,\widetilde V_c(w)\) est affine. La somme est donc strictement convexe sur le compact \(\Delta_K\), d'où l'existence et l'unicité du minimiseur.

Montrons maintenant que le minimiseur appartient à l'intérieur \(\Delta_K^\circ\). Supposons au contraire que \(p^\star_i=0\) pour un certain indice \(i\). Comme \(\sum_c p^\star_c=1\), il existe un indice \(j\) tel que \(p^\star_j>0\). Pour \(t\in(0,p^\star_j)\), posons \(p(t)=p^\star+t(e_i-e_j)\in\Delta_K\). On obtient une fonction unidimensionnelle dont la dérivée à droite en \(0\) est \(-\infty\), car \(\log t\to-\infty\) lorsque \(t\downarrow0\). Cela contredit la minimalité de \(p^\star\). Ainsi, tout minimiseur appartient à \(\Delta_K^\circ\).
\end{proof}

\begin{theorem}[Forme explicite du routage souple]
Le minimiseur unique de \(\Rcal_\beta(w,\cdot)\) est donné par
\[
p_c^\star(w)=\frac{\exp\!\bigl(-\beta\,\widetilde V_c(w)\bigr)}
{\sum_{j=1}^K \exp\!\bigl(-\beta\,\widetilde V_j(w)\bigr)}.
\]
\end{theorem}

\begin{proof}
Comme le minimiseur appartient à \(\Delta_K^\circ\), on peut appliquer la méthode des multiplicateurs de Lagrange au problème
\[
\min_{p\in\Delta_K^\circ}\ \sum_{c=1}^K p_c\,\widetilde V_c(w)+\frac{1}{\beta}\sum_{c=1}^K p_c\log p_c.
\]
On obtient, pour chaque \(c\),
\[
\widetilde V_c(w)+\frac{1}{\beta}(1+\log p_c)+\lambda=0,
\]
donc
\[
p_c=\exp(-1-\beta\lambda)\,\exp\!\bigl(-\beta\widetilde V_c(w)\bigr).
\]
La contrainte \(\sum_c p_c=1\) donne la formule annoncée. Le résultat est la forme usuelle du softmin entropique \cite{BauschkeCombettes2017,RockafellarWets1998}.
\end{proof}

\begin{corollary}[Valeur du problème de routage et limite dure]
On a l'identité
\[
\inf_{p\in\Delta_K}\Rcal_\beta(w,p)
=
-\frac{1}{\beta}\log\!\left(\sum_{c=1}^K \exp\!\bigl(-\beta\,\widetilde V_c(w)\bigr)\right).
\]
De plus,
\[
F(w)\le -\frac{1}{\beta}\log\!\left(\sum_{c=1}^K e^{-\beta \widetilde V_c(w)}\right)\le F(w)+\frac{\log K}{\beta}.
\]
Enfin, si le minimiseur de \(c\mapsto \widetilde V_c(w)\) est unique, noté \(c^\star(w)\), alors
\[
p^\star_\beta(w)\xrightarrow[\beta\to+\infty]{}e_{c^\star(w)}.
\]
\end{corollary}

\begin{proof}
C'est la valeur classique de l'entropie log-sum-exp obtenue en injectant la solution précédente dans \(\Rcal_\beta\). La borne sandwich est la propriété standard du \(\log\)-sum-exp. Pour la limite dure, on compare chaque poids non optimal au poids optimal ; les différences strictement positives dans les valeurs réduites sont écrasées exponentiellement lorsque \(\beta\to+\infty\).
\end{proof}

\begin{remark}
Le routage souple est une relaxation entropique du problème discret. Il ne remplace pas le routage exact ; il l'approxime de manière stable.
\end{remark}

% =================================================
\section{Pont avec PCH-ESM : base spectrale, spécialisation certifiée et apprentissage}
% =================================================

Cette section formalise le lien avec PCH-ESM de manière compatible et positive pour la théorie globale. L'idée n'est pas d'ajouter une architecture concurrente, mais d'identifier une instantiation canonique du jeu de contextes \(\Ccal\) par les blocs spécialisés certifiés fournis par PCH-ESM.

Le fichier PCH-ESM formalise trois briques pertinentes pour le routage :
(i) une décomposition orthogonale déterministe d'un latent en sous-espaces spécialisés ;
(ii) une fonctionnelle de spécialisation certifiée ;
(iii) une coordination globale par collage pondéré.
Dans ce cadre, le routage variationnel fournit la règle de sélection des branches ou blocs spécialisés.

\begin{definition}[Contextes PCH-ESM certifiés]
On note \(\Ccal_{\mathrm{ESM}}=\{1,\dots,n\}\) l'ensemble fini des blocs spécialisés obtenus par la décomposition spectrale de PCH-ESM. À chaque bloc \(i\in\Ccal_{\mathrm{ESM}}\) on associe une valeur réduite normalisée \(\widetilde V_i^{\mathrm{ESM}}(w)\) calculée par le problème local sur le bloc spécialisé correspondant.
\end{definition}

\begin{proposition}[Instantiation PCH-ESM du routage]
Si les branches locales du routage sont prises parmi les blocs certifiés \(\Ccal_{\mathrm{ESM}}\), alors le contexte choisi est
\[
i^\star(w)\in\argmin_{i\in\Ccal_{\mathrm{ESM}}}\widetilde V_i^{\mathrm{ESM}}(w).
\]
Autrement dit, la spécialisation certifiée de PCH-ESM fournit une famille de contextes admissibles, et le théorème de routage sélectionne le bloc de plus faible valeur réduite.
\end{proposition}

\begin{proof}
C'est une application directe du théorème de routage exact aux blocs spécialisés certifiés. Le fait que les blocs viennent d'une construction spectrale déterministe ne change pas la structure de l'argument : il fournit seulement une famille de contextes particulièrement bien contrôlée.
\end{proof}

\begin{figure}[ht]
\centering
\begin{tikzpicture}[node distance=1.6cm, >=Latex, every node/.style={align=center}]
\node (pch) [draw, rounded corners, minimum width=3.2cm, minimum height=0.95cm] {PCH\\valeurs locales};
\node (esm) [draw, rounded corners, right=1.8cm of pch, minimum width=3.5cm, minimum height=0.95cm] {PCH-ESM\\blocs certifiés};
\node (vals) [draw, rounded corners, right=1.8cm of esm, minimum width=3.4cm, minimum height=0.95cm] {Valeurs réduites\\par bloc};
\node (sel) [draw, rounded corners, below=1.2cm of vals, minimum width=3.4cm, minimum height=0.95cm] {Sélection\\\(\argmin\)};
\node (out) [draw, rounded corners, left=1.8cm of sel, minimum width=3.4cm, minimum height=0.95cm] {Contexte choisi\\ou bloc expert};

\draw[->] (pch) -- (esm);
\draw[->] (esm) -- (vals);
\draw[->] (vals) -- (sel);
\draw[->] (sel) -- (out);
\end{tikzpicture}
\caption{Chaîne compatissante entre PCH, PCH-ESM et routage : le bloc certifié fournit une branche locale, et le routage choisit le bloc de plus faible valeur réduite.}
\end{figure}

\begin{remark}
Cette compatibilité fait apparaître le vrai rôle de PCH-ESM dans la théorie globale : il ne remplace pas le routage, il fournit une famille de branches certifiées sur lesquelles le routage variationnel opère.
\end{remark}

% =================================================
\section{Théorème principal de synthèse}
% =================================================

\begin{theorem}[Structure variationnelle globale du routage PCH et son instantiation PCH-ESM]
Le routage PCH définit la structure variationnelle suivante :
\begin{enumerate}[leftmargin=2em]
\item pour chaque contexte \(c\), le problème local admet un unique minimiseur \(u_c^\star(w)=\prox_{\mu_c g_c}(w)\) et \(\widetilde V_c(w)=e_{\mu_c g_c}(w)\) ;
\item la valeur globale \(F(w)=\min_{c\in\Ccal}\widetilde V_c(w)\) est continue, localement lipschitzienne, directionnellement différentiable, et de classe \(C^1\) sur toute cellule de dominance stricte ;
\item l'espace latent se décompose en cellules strictes \(\Omega_c^{\mathrm{str}}\), en cellules faibles \(\Omega_c\), et en une strate de bord \(\Bcal\), et le contexte routé est localement constant sur toute cellule stricte ;
\item le routage souple est la régularisation entropique du critère des valeurs réduites, avec minimiseur de Gibbs ;
\item la correspondance active est semi-continue supérieurement, conformément au cas discret du théorème du maximum de Berge ;
\item le sous-différentiel de Clarke de \(F\) est l'enveloppe convexe des gradients actifs ;
\item lorsque la famille de contextes est fournie par les blocs spécialisés certifiés de PCH-ESM, le bloc routé est exactement le bloc spécialisé minimisant la valeur réduite normalisée.
\end{enumerate}
\end{theorem}

\begin{proof}
Les points \(1\) à \(6\) résultent des sections précédentes sur la normalisation, la valeur globale, la partition de l'espace latent, la relaxation entropique, la stabilité sous marge et la régularité de Clarke. Le point \(7\) est l'instantiation de la théorie abstraite à la famille de contextes certifiés de PCH-ESM. Le théorème résume donc des résultats déjà démontrés, mais il ajoute une conclusion structurelle nouvelle : la version PCH-ESM du routage n'est rien d'autre qu'une application du même principe aux blocs spectraux certifiés.
\end{proof}

% =================================================
\section{Exemple illustratif en dimension un}
% =================================================

\begin{example}[Deux contextes symétriques]
On prend \(d=1\), \(\Mcal_1=\{-a\}\), \(\Mcal_2=\{a\}\) avec \(a>0\), et on suppose les mêmes paramètres pour les deux contextes :
\[
\lambda_{\mathrm{ctx},1}=\lambda_{\mathrm{ctx},2}=\lambda>0,
\qquad
\lambda_{\mathrm{ridge},1}=\lambda_{\mathrm{ridge},2}=\rho\ge 0.
\]
Alors \(\mu_1=\mu_2=1+2\rho\), et une complétion du carré donne
\[
u_1^\star(w)=\frac{w-2\lambda a}{1+2\lambda+2\rho},
\qquad
u_2^\star(w)=\frac{w+2\lambda a}{1+2\lambda+2\rho}.
\]
De plus, la différence des valeurs réduites normalisées est affine :
\[
\widetilde V_1(w)-\widetilde V_2(w)
=
\frac{4a\lambda}{(1+2\rho)(1+2\lambda+2\rho)}\,w.
\]
Par conséquent,
\[
\Omega_1^{\mathrm{str}}=(-\infty,0),
\qquad
\Omega_2^{\mathrm{str}}=(0,\infty),
\qquad
\Bcal=\{0\}.
\]
Le routage exact se lit donc immédiatement : le contexte \(1\) est choisi pour \(w<0\), le contexte \(2\) pour \(w>0\), et la frontière de décision est ici réduite au point \(0\).
\end{example}

\begin{remark}
Cet exemple montre que les cellules ne sont pas un objet abstrait : elles décrivent une partition géométrique explicite de l'espace latent.
\end{remark}

% =================================================
\section{Relation à la littérature et positionnement}
% =================================================

Le cadre ici présenté se place naturellement à l'interface de l'analyse convexe, de la théorie des fonctions valeur, de la programmation disjonctive et des modèles de type \emph{mixture of experts}. Les références classiques pertinentes sont celles de Moreau, Rockafellar, Berge, Clarke, ainsi que Bauschke--Combettes et Rockafellar--Wets \cite{Moreau1965,Rockafellar1970,Berge1963,Clarke1983,BauschkeCombettes2017,RockafellarWets1998}.

La nouveauté du manuscrit n'est pas d'inventer de nouveaux outils de convexité ; elle est de faire du routage un problème variationnel à part entière. Les méthodes classiques de routage utilisent un score appris, un classifieur ou un gate séparé. Ici, le contexte est sélectionné comme minimiseur d'une famille de valeurs réduites issues directement de problèmes locaux. Le routage n'est donc plus un module externe : il devient une conséquence de la structure énergétique.

En termes pratiques, le routage appris est souvent moins coûteux à l'inférence une fois le gate entraîné, mais il introduit une phase d'apprentissage séparée, une dépendance aux données d'entraînement du gate, et une opacité de décision. Le routage variationnel est plus coûteux si l'on résout exhaustivement toutes les branches, mais il est certifiable, il s'inscrit directement dans PCH, et il peut être accéléré par un pré-roulage court avec certificat de marge. Ce compromis est précisément la bonne formulation scientifique du problème.

Autrement dit, la quantité fondamentale n'est pas un score appris séparément, mais l'objet
\[
\min_{c\in\Ccal} e_{\mu_c g_c}(w).
\]
C'est cette enveloppe inférieure qui porte l'information de décision, et non une sortie de classifieur externe. En ce sens, la géométrie de décision est induite par les branches locales elles-mêmes.

\begin{figure}[ht]
\centering
\begin{tikzpicture}[node distance=1.45cm, >=Latex, every node/.style={align=center}]
\node (loc) [draw, rounded corners, minimum width=3.8cm, minimum height=0.9cm] {Problèmes locaux\\$\widetilde J_c(\cdot;w)$};
\node (val) [draw, rounded corners, below=of loc, minimum width=3.6cm, minimum height=0.9cm] {Fonctions valeur\\$\widetilde V_c(w)$};
\node (inf) [draw, rounded corners, below=of val, minimum width=3.8cm, minimum height=0.9cm] {Enveloppe inférieure\\$F(w)=\min_c \widetilde V_c(w)$};
\node (geo) [draw, rounded corners, below=of inf, minimum width=4.0cm, minimum height=0.9cm] {Géométrie des cellules\\et des frontières};
\node (route) [draw, rounded corners, below=of geo, minimum width=3.6cm, minimum height=0.9cm] {Routage choisi\\$c^\star(w)$};

\draw[->] (loc) -- (val);
\draw[->] (val) -- (inf);
\draw[->] (inf) -- (geo);
\draw[->] (geo) -- (route);
\end{tikzpicture}
\caption{Chaîne conceptuelle de l'apport du manuscrit : du problème local à la géométrie de décision.}
\end{figure}

\begin{remark}
Cette lecture permet de séparer proprement ce qui relève d'un rappel classique (Moreau, Clarke, Berge, Rockafellar--Wets) et ce qui relève de la contribution propre du manuscrit : la sélection de contexte par minimisation des valeurs réduites et l'interprétation géométrique qui en découle.
\end{remark}

% =================================================
\section{Conclusion}
% =================================================

Le noyau PCH n'est pas le problème à résoudre ; il est déjà la brique mathématique centrale. Le vrai problème est le routage : comment choisir le contexte de manière rigoureuse, sans introduire un contournement externe.

La réponse proposée ici est écrite de façon exacte : le routage est un problème d'optimisation sur l'espace produit \(\R^d\times\Ccal\), et sa relaxation entropique sur le simplexe fournit une version souple. La formulation canonique est
\[
(\hat u(w),\hat c(w))\in\argmin_{(u,c)\in\R^d\times\Ccal}\mathcal E_{\mathrm{route}}(w;u,c),
\]
équivalente à
\[
\hat c(w)\in\argmin_{c\in\Ccal}\widetilde V_c(w),
\qquad
\widetilde V_c(w)=e_{\mu_c g_c}(w).
\]
La version souple est donnée par
\[
p^\star_\beta(w)=\argmin_{p\in\Delta_K}\Rcal_\beta(w,p).
\]

Cette formulation conserve strictement le noyau PCH, définit explicitement le critère de sélection du contexte, donne une écriture exacte du routage comme problème sur \(\R^d\times\Ccal\), admet une relaxation convexe sur le simplexe, relie le pré-roulage à la convergence vers la valeur réduite, et fournit une interprétation variationnelle propre du choix de contexte. Le pont avec PCH-ESM est lui aussi explicite : la décomposition spectrale certifiée fournit une famille de branches finies sur lesquelles le même principe de minimisation sélectionne le bloc spécialisé.

\begin{thebibliography}{99}

\bibitem{Moreau1965}
J.-J. Moreau,
\newblock Proximité et dualité dans un espace hilbertien.
\newblock \emph{Bull. Soc. Math. France}, 93:273--299, 1965.

\bibitem{Rockafellar1970}
R. T. Rockafellar,
\newblock \emph{Convex Analysis}.
\newblock Princeton University Press, 1970.

\bibitem{Berge1963}
C. Berge,
\newblock \emph{Topological Spaces: Including a Treatment of Multi-Valued Functions, Vector Spaces and Convexity}.
\newblock Oliver \& Boyd, 1963.

\bibitem{Clarke1983}
F. H. Clarke,
\newblock \emph{Optimization and Nonsmooth Analysis}.
\newblock Wiley, 1983.

\bibitem{BauschkeCombettes2017}
H. H. Bauschke and P. L. Combettes,
\newblock \emph{Convex Analysis and Monotone Operator Theory in Hilbert Spaces}.
\newblock 2nd ed., Springer, 2017.

\bibitem{RockafellarWets1998}
R. T. Rockafellar and R. J.-B. Wets,
\newblock \emph{Variational Analysis}.
\newblock Springer, 1998.

\end{thebibliography}

\end{document}
